Das Zimmer ist nicht dicht, der Druck bleibt gleich dem äusseren Luftdruck. Die Luft, die am Morgen im Zimmer war, nimmt bei der höheren Temperatur nach Gay-Lussac das Volumen
\begin{align}
 V_2 &= \VzZF = \Vz\cdot\frac{\TZw}{\TEi} = \VzZ
\end{align}
ein, mit $T_1 = \TEiF = \TEi$ und $T_2 = \TZwF = \TZw$. Das Zimmer fasst aber nur $V_1$; der Rest entweicht:
\begin{align}
 \Delta V &= \dVzF = \dVz \approx \result{\dVzP{0}{2}}
\end{align}
